
\documentclass{article}
\usepackage{colortbl}
\usepackage{makecell}
\usepackage{multirow}
\usepackage{supertabular}

\begin{document}

\newcounter{utterance}

\centering \large Interaction Transcript for game `cladder', experiment `full\_v1.5\_default', episode 525 with refineonly{-}Qwen3.5{-}4B.
\vspace{24pt}

{ \footnotesize  \setcounter{utterance}{1}
\setlength{\tabcolsep}{0pt}
\begin{supertabular}{c@{$\;$}|p{.15\linewidth}@{}p{.15\linewidth}p{.15\linewidth}p{.15\linewidth}p{.15\linewidth}p{.15\linewidth}}
    \# & \multicolumn{2}{c}{Player} && \multicolumn{2}{c}{Game Master} \\
    \hline

    \theutterance \stepcounter{utterance}  
    & & & \multicolumn{4}{p{0.6\linewidth}}{
        \cellcolor[rgb]{0.9,0.9,0.9}{
            \makecell[{{p{\linewidth}}}]{
                \texttt{\tiny{[P1$\langle$GM]}}
                \texttt{You are an expert in causal inference. The following question is not a typical commonsense query, but rather a meticulously designed question created by a professor specializing in causal inference, intended to assess the students' mastery of the course content.} \\
\texttt{Imagine a self{-}contained, hypothetical world with only the following conditions, and without any unmentioned factors or causal relationships: Unobserved confounders has a direct effect on drug taken and cholesterol level. Treatment assignment has a direct effect on drug taken. Drug taken has a direct effect on cholesterol level. Unobserved confounders is unobserved. For patients not assigned the drug treatment, the probability of low cholesterol is 63\%. For patients assigned the drug treatment, the probability of low cholesterol is 52\%. For patients not assigned the drug treatment, the probability of taking of all assigned drugs is 37\%. For patients assigned the drug treatment, the probability of taking of all assigned drugs is 74\%. Will taking of all assigned drugs increase the chance of low cholesterol?} \\
\texttt{Start your answer with 'Yes' or 'No', followed by additional reasoning or evidence to support your explanation.} \\
            }
        }
    }
     \\ \\

    \theutterance \stepcounter{utterance}  
    & & \multicolumn{4}{p{0.6\linewidth}}{
        \cellcolor[rgb]{0.95,0.95,0.95}{
            \makecell[{{p{\linewidth}}}]{
                \texttt{\tiny{[P1$\leftrightarrow$P]}}
                \texttt{{[}DRAFT{]}} \\
\texttt{No} \\
\\ 
\texttt{**Reasoning:**} \\
\\ 
\texttt{To determine whether taking the drug increases the chance of low cholesterol, we must calculate the **Average Treatment Effect on the Treated (ATT)**, which represents the causal effect of the drug on the subgroup of patients who actually took it. The ATT is calculated as the difference between the potential outcome of low cholesterol for those who took the drug (\$Y\_\{1\}\$) and the potential outcome for those who did not take the drug (\$Y\_\{0\}\$), restricted to the population that took the drug.} \\
\\ 
\texttt{The formula for the ATT is:} \\
\texttt{\$\$ATT = P(Y\_\{1\}=1 | \textbackslash{}text\{took drug\}) {-} P(Y\_\{0\}=1 | \textbackslash{}text\{took drug\})\$\$} \\
\\ 
\texttt{We are given the following data points:} \\
\texttt{1.  **\$P(Y\_\{1\}=1 | \textbackslash{}text\{took drug\})\$**: The probability of low cholesterol for patients assigned the drug who took it is **52\%**.} \\
\texttt{2.  **\$P(Y\_\{0\}=1 | \textbackslash{}text\{took drug\})\$**: The probability of low cholesterol for patients assigned the drug who *did not* take it.} \\
\\ 
\texttt{To find the second value, we use the **Law of Total Probability**. We know the overall probability of low cholesterol for patients assigned the drug (\$P(Y\_\{1\}=1 | \textbackslash{}text\{assigned\})\$) is 52\%. This overall probability is a weighted average of the outcome for those who took the drug and those who did not, based on the probability of taking the drug within that group.} \\
\\ 
\texttt{*   Probability of taking the drug given assignment: \$P(\textbackslash{}text\{took\} | \textbackslash{}text\{assigned\}) = 74\textbackslash{}\% = 0.74\$.} \\
\texttt{*   Probability of not taking the drug given assignment: \$P(\textbackslash{}text\{not took\} | \textbackslash{}text\{assigned\}) = 1 {-} 0.74 = 0.26\$.} \\
\texttt{*   Overall probability of low cholesterol given assignment: \$P(Y\_\{1\}=1 | \textbackslash{}text\{assigned\}) = 52\textbackslash{}\% = 0.52\$.} \\
\\ 
\texttt{The equation is:} \\
\texttt{\$\$P(Y\_\{1\}=1 | \textbackslash{}text\{assigned\}) = P(Y\_\{1\}=1 | \textbackslash{}text\{took\}, \textbackslash{}text\{assigned\}) \textbackslash{}cdot P(\textbackslash{}text\{took\} | \textbackslash{}text\{assigned\}) + P(Y\_\{1\}=1 | \textbackslash{}text\{not took\}, \textbackslash{}text\{assigned\}) \textbackslash{}cdot P(\textbackslash{}text\{not took\} | \textbackslash{}text\{assigned\})\$\$} \\
\\ 
\texttt{Substituting the known values:} \\
\texttt{\$\$0.52 = (0.52 \textbackslash{}cdot 0.74) + P(Y\_\{0\}=1 | \textbackslash{}text\{took drug\}) \textbackslash{}cdot 0.26\$\$} \\
\\ 
\texttt{Solving for \$P(Y\_\{0\}=1 | \textbackslash{}text\{took drug\})\$:} \\
\texttt{\$\$0.52 = 0.38} \\
            }
        }
    }
    & \\ \\

    \theutterance \stepcounter{utterance}  
    & \multicolumn{4}{p{0.6\linewidth}}{
        \cellcolor[rgb]{0.561,0.737,0.561}{
            \makecell[{{p{\linewidth}}}]{
                \texttt{\tiny{[P1$\rangle$GM]}}
                \texttt{No} \\
\\ 
\texttt{**Reasoning:**} \\
\\ 
\texttt{To determine whether taking the drug increases the chance of low cholesterol, we must calculate the **Average Treatment Effect on the Treated (ATT)**, which represents the causal effect of the drug on the subgroup of patients who actually took it. The ATT is calculated as the difference between the potential outcome of low cholesterol for those who took the drug (\$Y\_\{1\}\$) and the potential outcome for those who did not take the drug (\$Y\_\{0\}\$), restricted to the population that took the drug.} \\
\\ 
\texttt{The formula for the ATT is:} \\
\texttt{\$\$ATT = P(Y\_\{1\}=1 | \textbackslash{}text\{took drug\}) {-} P(Y\_\{0\}=1 | \textbackslash{}text\{took drug\})\$\$} \\
\\ 
\texttt{We are given the following data points:} \\
\texttt{1.  **\$P(Y\_\{1\}=1 | \textbackslash{}text\{took drug\})\$**: The probability of low cholesterol for patients assigned the drug who took it is **52\%**.} \\
\texttt{2.  **\$P(Y\_\{0\}=1 | \textbackslash{}text\{took drug\})\$**: The probability of low cholesterol for patients assigned the drug who *did not* take it.} \\
\\ 
\texttt{To find the second value, we use the **Law of Total Probability**. We know the overall probability of low cholesterol for patients assigned the drug (\$P(Y\_\{1\}=1 | \textbackslash{}text\{assigned\})\$) is 52\%. This overall probability is a weighted average of the outcome for those who took the drug and those who did not, based on the probability of taking the drug within that group.} \\
\\ 
\texttt{*   Probability of taking the drug given assignment: \$P(\textbackslash{}text\{took\} | \textbackslash{}text\{assigned\}) = 74\textbackslash{}\% = 0.74\$.} \\
\texttt{*   Probability of not taking the drug given assignment: \$P(\textbackslash{}text\{not took\} | \textbackslash{}text\{assigned\}) = 1 {-} 0.74 = 0.26\$.} \\
\texttt{*   Overall probability of low cholesterol given assignment: \$P(Y\_\{1\}=1 | \textbackslash{}text\{assigned\}) = 52\textbackslash{}\% = 0.52\$.} \\
\\ 
\texttt{The equation is:} \\
\texttt{\$\$P(Y\_\{1\}=1 | \textbackslash{}text\{assigned\}) = P(Y\_\{1\}=1 | \textbackslash{}text\{took\}, \textbackslash{}text\{assigned\}) \textbackslash{}cdot P(\textbackslash{}text\{took\} | \textbackslash{}text\{assigned\}) + P(Y\_\{1\}=1 | \textbackslash{}text\{not took\}, \textbackslash{}text\{assigned\}) \textbackslash{}cdot P(\textbackslash{}text\{not took\} | \textbackslash{}text\{assigned\})\$\$} \\
\\ 
\texttt{Substituting the known values:} \\
\texttt{\$\$0.52 = (0.52 \textbackslash{}cdot 0.74) + P(Y\_\{0\}=1 | \textbackslash{}text\{took drug\}) \textbackslash{}cdot 0.26\$\$} \\
\\ 
\texttt{Solving for \$P(Y\_\{0\}=1 | \textbackslash{}text\{took drug\})\$:} \\
\texttt{\$\$0.52 = 0.38} \\
            }
        }
    }
    & & \\ \\

    \theutterance \stepcounter{utterance}  
    & & & \multicolumn{2}{p{0.3\linewidth}}{
        \cellcolor[rgb]{0.9,0.9,0.9}{
            \makecell[{{p{\linewidth}}}]{
                \texttt{\tiny{[GM$|$GM]}}
                \texttt{No} \\
            }
        }
    }
    & & \\ \\

    \theutterance \stepcounter{utterance}  
    & & & \multicolumn{2}{p{0.3\linewidth}}{
        \cellcolor[rgb]{0.9,0.9,0.9}{
            \makecell[{{p{\linewidth}}}]{
                \texttt{\tiny{[GM$|$GM]}}
                \texttt{no} \\
            }
        }
    }
    & & \\ \\

    \theutterance \stepcounter{utterance}  
    & & & \multicolumn{2}{p{0.3\linewidth}}{
        \cellcolor[rgb]{0.9,0.9,0.9}{
            \makecell[{{p{\linewidth}}}]{
                \texttt{\tiny{[GM$|$GM]}}
                \texttt{game\_result = WIN} \\
            }
        }
    }
    & & \\ \\

\end{supertabular}
}

\end{document}
